\section{Limitations}

The main constraints of this study are the limited amount of data and the lack of detailed trip information. The first dataset has an 8-hour resolution, which complicates accurate model training and evaluation. Although the second dataset offers a 30-minute resolution, it is limited to two weeks and does not include bikes in transit.
In addition to the battery data, the incomplete trip data further complicates model development. Information on the exact routes taken by users is unavailable, making it difficult to estimate battery consumption for round trips. Additionally, details such as the bike's condition (e.g., tire inflation) or relevant rider characteristics (e.g., weight or height) are also missing.

\section{Discussion}

The analysis revealed distinct spatiotemporal variations in \ac{e-b} battery levels. Battery levels were generally lower in the evenings, reflecting the impact of daily usage, and higher in the mornings after overnight resting periods. Stations in northern neighborhoods exhibited higher average battery levels, likely attributed to lower usage rates and longer resting times. Additionally, notable correlations were observed between stations' battery levels, their inter-event times, and their distance from other stations.

Moreover, the developed distance index provides valuable insights into one of the primary operational challenges faced by electric \acp{bss}: battery level management. Stations with lower distance index values are more prone to accelerated battery depletion, as they tend to receive more frequent trips from nearby origins, leaving less time for recharging between trips. This metric offers \ac{bss} operators a practical tool to identify such stations and prioritize interventions, including more frequent battery replacements, recharging, or strategic bike redistribution, thereby enhancing the system's efficiency and reliability.

The presented bike battery forecasting methodology combines a traditional probabilistic framework with heterogeneous data sources, resulting in a multifaceted tool capable of improving the real-time management of micro-mobility services. Although some challenges remain due to the lack of on-route data, such as speed or location, the battery function has been validated. This function, when combined with a Markov chain approach, provides a validated prediction for the availability of \acp{e-b}, either due to an imbalance between stations, or to an overused battery. Further improvements could be made, particularly if more detailed data were available. Access to \ac{gps} information on the routes taken would likely improve power consumption estimates, while higher-resolution data could enhance battery predictions through machine learning approaches. Additionally, experimental data collected through ad-hoc studies could improve the model's performance by fine-tuning the parameters, ensuring a more accurate representation of real-world conditions.

With the progressive electrification of bike fleets worldwide, predictions of availability alone may be insufficient, as docked bikes with low batteries remain unusable. The presented approach opens the door to developing microscopic models that provide joint predictions of bike availability and battery levels. This framework could support \ac{bss} providers in optimizing bike availability and identifying stations at higher risk of disrupting urban mobility.